%!TEX root = ../main.tex
\section{Cost-Aware Execution}
\label{sec:executor}

A memory system that answers well but spends unboundedly does not scale to long histories and heavy usage. This section prices what \sys{} spends and then shows how to spend less of it without giving up accuracy.

\stitle{Accounting.}
\sys{} issues LLM calls on both the write path (catalog extraction) and the read path (query reasoning). Write-path cost is paid once per session and amortizes across queries. The rest of this section concerns read-path cost, which dominates. The standard remedy for accuracy, self-consistency at $K{=}3$~\cite{sc}, triples that cost on \emph{every} query, regardless of whether the extra samples change the answer. Formally, the total cost over a workload of $Q$ queries with class fractions $(p_{\point},p_{\agg},p_{\cmp})$ under fixed budgets $(K_{\point},K_{\agg},K_{\cmp})$ is
\begin{equation}
\label{eq:cost}
C \;=\; Q\cdot\bigl(p_{\point}K_{\point} + p_{\agg}K_{\agg} + p_{\cmp}K_{\cmp}\bigr).
\end{equation}
A Uniform-$K{=}3$ reference sets $K_{\point}=K_{\agg}=K_{\cmp}=3$, so $C=3Q$.

The executor makes this budget an explicit, plan-aware decision, exploiting two independent sources of waste. One is the class of queries saturated at a single call (Section~\ref{subsec:budget}). The other is the individual query whose answer is stable before its budget is spent (Section~\ref{subsec:shortcircuit}). The two act on disjoint slices, and their savings compose.

\subsection{Plan-Class Budgeting}
\label{subsec:budget}

The first source of waste is structural, since one class of query never benefits from the samples it is given. This subsection assigns a budget per class and derives what that assignment saves as a function of workload composition. The executor pre-declares a per-class budget instead of a per-query one, setting $K_{\point}=1$ and $K_{\agg}=K_{\cmp}=3$. The asymmetry follows from what each plan computes, not from a heuristic. A \point{} query resolves a single stored fact. Once $\textsf{(Alice, has-provider, Raiffeisen Basel)}$ is in the grounding block, replicated sampling can only re-read it, and a single call is both necessary and sufficient. \agg{} and \cmp{} queries compose evidence across multiple records, as when the outdoor count must decide eight independent \texttt{yes}/\texttt{no} judgments, and there sampling genuinely reduces variance, so these classes keep the full budget. Substituting these budgets into Eq.~(\ref{eq:cost}), the saving against Uniform-$K{=}3$ is $\Delta C=2p_{\point}Q$. On \membench{} with $p_{\point}=0.61$, this alone reduces $C$ from $900$ to $532$, a $40.9\%$ cut. The saving grows with the share of point queries and costs nothing in accuracy on exactly the queries it economizes, in line with the DBMS discipline of pricing an operator by what it must compute and not by a uniform per-tuple rate.

\subsection{Stability Short-Circuit}
\label{subsec:shortcircuit}

The second source of waste lies inside the hard classes, where some instances are easy but indistinguishable from the rest before execution. This subsection recovers that distinction at no extra cost by reading the agreement between samples the system was going to draw anyway. Within the \agg/\cmp{} budget, the executor decides adaptively whether the third call is worth making. It draws two independent samples $p_1,p_2$ under the operator's template, measures their token-level agreement $\mathrm{agree}(\cdot,\cdot)\in[0,1]$ against a threshold $\theta$, and draws a third sample $p_3$ only when the two disagree, giving the rule
\begin{equation}
\label{eq:short}
\hat{a} \;=\; \begin{cases}
p_1, & \mathrm{agree}(p_1,p_2)\geq\theta, \\
\mathrm{Majority}(p_1,p_2,p_3), & \text{otherwise.}
\end{cases}
\end{equation}
Under Eq.~(\ref{eq:short}) on Alice's outdoor question, both samples enumerate the same eight events and emit \texttt{Answer: 3}, so the agreement is high, the third call is skipped, and the query costs two. Had one sample returned \texttt{3} and the other \texttt{4}, the low agreement would have triggered escalation, spending the third call where it can change the outcome. We fix $\theta=0.6$ for all benchmarks and all backbones, and it is never retuned per dataset. Combined with plan-class budgeting, agreement-based stopping drives the total cost on \membench{} from $532$ to $444$, a further $16.5\%$ cut and $50.7\%$ overall against the Uniform-$K{=}3$ reference.

\stitle{Agreement as a Cost Signal.}
The reinterpretation here is the contribution. Self-consistency~\cite{sc} samples $K$ reasoning paths and votes, treating disagreement as \emph{noise to average away} under a sampling budget fixed in advance. \sys{} reads the same disagreement as \emph{information about the query}. When two samples agree, the query is easy for the current model and grounding, and a third sample is unlikely to change the outcome. Agreement thus becomes a stopping rule, not a voting mechanism. Adaptive-Consistency~\cite{adaptiveconsistency} and early-stopping self-consistency~\cite{escsc} both halt sampling once drawn answers converge, but they operate on a single query in isolation and must draw at least two samples before learning anything. \sys{} knows from Eq.~(\ref{eq:classifier}) that a point lookup cannot benefit from a second sample at all, and it skips the extra sample entirely for the majority of the workload. Agreement-based stopping and class-based budgeting are therefore complementary, the first acting within a class and the second across classes, and Section~\ref{subsec:exp4} measures them separately.
